\section*{Broader Impact Statement}
\label{sec:broader_impact}

\paragraph{Positive Impacts.}
Efficient movie scene scheduling directly reduces production costs, which can make filmmaking more accessible to independent producers and lower-budget projects.
By optimizing actor utilization and minimizing idle days, our approach reduces waste in the form of unnecessary travel, idle crew time, and extended shooting schedules.
The explicit modeling of actor consecutive working-day limits (constraint~\ref{c13}) promotes compliance with labor regulations and union rules, ensuring that generated schedules respect worker welfare by construction rather than relying on post-hoc manual correction.
The three-tier location hierarchy and town-continuity constraint (\ref{c12}) produce schedules with coherent geographic progression, reducing disruptive overnight relocations for cast and crew.
The decomposition methodology transfers to other resource-constrained scheduling domains including construction project management, operating room scheduling, and manufacturing planning.

\paragraph{Potential Concerns.}
Automated scheduling tools could reduce human involvement in creative production decisions.
Film scheduling involves nuanced factors---actor fatigue, creative flow, weather dependencies---that our cost-centric optimization does not fully capture, despite the criticality weights and labor constraints providing partial modeling of these concerns.
We emphasize that BDM-AOS is intended as a decision-support tool that provides optimized candidate schedules for human producers to review and modify, not as a replacement for human judgment.

\paragraph{Environmental Considerations.}
Optimized scheduling with hierarchical transfer cost minimization directly reduces unnecessary inter-town and inter-city location transfers, which decreases transportation-related carbon emissions.
The town-continuity constraint (\ref{c12}) further reduces emissions by eliminating back-and-forth travel between towns on consecutive days.
The computational cost of running BDM-AOS (typically under 10 seconds per instance) is negligible compared to the environmental and financial savings from even marginal schedule improvements.
